\documentclass[12pt]{article}
\usepackage[utf8]{inputenc}
\usepackage[T1]{fontenc}
\usepackage{hyperref}
\title{Tawallammat Tamajaq Language Resources for K-12 Curriculum Development}
\author{Compiled by Claude Sonnet 5 (claude-sonnet-5) for Andr\'as Kornai}
\date{}
\begin{document}
\maketitle

\section{Introduction}
Tawallammat Tamajaq (ISO 639-3: ttq; Glottocode tawa1286) is a Southern Tamasheq variety of the Berber branch of Afro-Asiatic, spoken mainly in Niger and Mali. Speaker estimates differ: Omniglot gives about 870,000, and Wikipedia gives about 1.3 million for 2021--2022. Omniglot says it is written in Latin, Arabic and Tifinagh-based script. It also says it is taught in some primary schools and used in literature and on the radio. I could not confirm that claim against a second source.

\section{Primary Educational Resources}
No primers or textbooks were found in the searches. Omniglot's statement about some primary-school use (above) is the only school-related information found.
\begin{itemize}
\item \url{https://omniglot.com/writing/tawallammattamajaq.htm}
\end{itemize}

\section{Linguistic and Reference Materials}
\begin{itemize}
\item Wikipedia entry for Tawellemet: \url{https://wikipedia.com/wiki/ISO_639:ttq}
\item Enguehard and Modi (2009), ``Towards an Electronic Dictionary of Tamajaq Language in Niger''. This is about Tamajaq in Niger generally and may not cover this variety specifically: \url{https://aflat.org/node/316}
\item ACL Anthology copy of the same paper (PDF): \url{https://preview.aclanthology.org/moar-dois/W09-0712.pdf}
\end{itemize}

\section{Religious and Community Resources}
\begin{itemize}
\item ScriptureEarth page for ttq: \url{https://www.scriptureearth.org/00i-Scripture_Index.php?iso=ttq}
\item Words of Life audio recordings, Tamajeq East Tawallemet: \url{https://megavoice.com/media-cloud/m025853-words-of-life-ttq-tamajeq-east-tawallemet-audio-bible/}
\end{itemize}

\section{Digital and Audio Resources}
\begin{itemize}
\item Global Recordings Network, Tamajeq (Tamastrait; a related variety, not necessarily ttq): \url{https://www.globalrecordings.net/en/language/1793}
\end{itemize}

\section{Curriculum Development Context}
This is the best resourced of the four languages, with a large and stable speaker base, a documented orthography and scripture material. Textbooks and curriculum documents were not found in these searches, and the extent of school use still needs checking with Nigerien and Malian education authorities and local organisations.
\end{document}